\scanpage{10}
\mspage{4}
Le graphe indique comment les composantes irréductibles de $X$ (figurés par les sommets) se recoupent deux à deux -- les pts d'intersection, i.e. les pts singuliers (\inserted{``doubles'', \unsure{sc.|ie.}} quadratiques) de $X$, correspondent aux arêtes. Si une composante irréductible $X_\alpha$ de $X$ correspond au sommet $\alpha$ du graphe, alors les boucles fermées en $\alpha$ correspondent biunivoquement aux pts doubles de $X_\alpha$ -- donc dire que les $X_\alpha$ sont lisses signifie que le graphe $G_X$ n'a pas de boucle fermée, i.e. que pour \unsure{tt|tte} arc, l'origine est $\neq$ de l'extrémité.

Il est clair que tout graphe fini peut être obtenu \inserted{(à iso près)} \deleted{\unsure{à} l'aide} par une $\widehat X$ convenable -- et $=$ avec des \unsure{compo-|comp.-}
\mspage{9}
